\chapter{Bulk dump}

Bulk dump transfers the instrument's internal data as System Exclusive. The instrument is
silent for the duration, and the transfer uses the MIDI~1 sockets only.

\section{Categories and what each transmits}
The \textsc{sysex bulk dump} screen offers five categories. Each sends one or more data
messages at fixed dump addresses.

\begin{center}
\begin{tabular}{llllr}
\toprule
category & address bytes & length bytes & length & blocks \\
\midrule
\textsc{system, part \& midi} & \bytes{40 00 00} & \bytes{00 00 20} & 32 & \\
                              & \bytes{40 00 20} & \bytes{00 12 60} & 2\,400 & \\
\addlinespace
\textsc{sound}                & \bytes{20 00 00} & \bytes{10 00 00} & 262\,144 & 32 \\
\addlinespace
\textsc{combination}          & \bytes{50 00 00} & \bytes{00 06 00} & 768 & \\
                              & \bytes{50 06 00} & \bytes{05 40 00} & 90\,112 & 16 \\
\addlinespace
\textsc{sequencer}            & \bytes{60 00 00} & \bytes{00 18 00} & 3\,072 & \\
                              & \bytes{60 18 00} & \bytes{01 70 00} & 30\,720 & \\
                              & \bytes{62 08 00} & \multicolumn{2}{l}{set at send time} & \\
\addlinespace
\textsc{total keyboard}       & \multicolumn{4}{l}{the four categories in the order above} \\
\bottomrule
\end{tabular}
\end{center}

Lengths are in source bytes; on the wire each occupies twice as many, because of the
payload encoding described below.

The \emph{blocks} column is internal structure, not extra messages: \textsc{sound} sends
one header covering all 262\,144 bytes, which is 32 repeats of an 8\,192-byte layout, and
\textsc{combination}'s second block one header covering 90\,112 bytes as 16 repeats of
5\,632. Those are the instrument's 32 sounds and 16 combinations.

\textsc{total keyboard} runs the four categories one after another in the order listed,
each closed by its own end-of-category message, with a single end-of-dump message after
the last.

The \textsc{sequencer} category is the only one whose final message has no fixed length ---
it carries as much recorded material as the instrument holds.

\begin{note}{The dump address is not a memory address}
Within one category the addresses are consecutive: a block's address is the previous
block's address plus its length. Across categories they are not. \textsc{sound} and
\textsc{combination} are held side by side inside the instrument, yet their dump addresses
are 786\,432 apart; other pairs are the reverse.

Treat the address and length bytes together as a fixed name for one of the eight blocks.
The instrument does the same --- on reception it compares them against a fixed list and
performs no arithmetic on them.
\end{note}

\section{Data message}
\begin{center}
\bytes{F0 50 2D 04 \textit{nn} 11}\quad
\bytes{\textit{a2 a1 a0}}\quad
\bytes{\textit{l2 l1 l0}}\quad
\bytes{\textit{data\ldots}}\quad
\bytes{\textit{c}}\quad
\bytes{\textit{sum}}\quad
\bytes{F7}
\end{center}

\begin{description}
\item[\bytes{\textit{nn}}] unit field, \bytes{00} or \bytes{01}.
\item[\bytes{\textit{a2 a1 a0}}] destination address, three bytes, most significant first,
  seven bits each: the value is $a_2 \times 128^2 + a_1 \times 128 + a_0$.
\item[\bytes{\textit{l2 l1 l0}}] length in \emph{source} bytes, encoded the same way.
\item[\bytes{\textit{data}}] the payload, encoded as described below.
\item[\bytes{\textit{c}}] continuation flag: \bytes{01} if further messages follow,
  \bytes{00} on the last.
\item[\bytes{\textit{sum}}] checksum.
\end{description}

A message is filled to at most 252 bytes before being sent; longer transfers are split.

\section{Payload encoding}
Each source byte is sent as two message bytes, \textbf{high nibble first}:
\[ b \;\longrightarrow\; (b \gg 4),\quad (b \mathbin{\&} \mathtt{0F}) \]
So source byte \bytes{F1} appears on the wire as \bytes{0F 01}, and the payload occupies
twice as many bytes as the length field states. A receiver rejoins them as
$((\text{first} \mathbin{\&} \mathtt{0F}) \ll 4) \mathbin{|} (\text{second} \mathbin{\&}
\mathtt{0F})$.

An odd number of payload bytes is invalid.

\section{Checksum}
Let $S$ be the eight-bit sum of every byte from the identifier \bytes{50} through the
continuation flag inclusive. Then
\[ \mathit{sum} \;=\; (0 - S) \mathbin{\&} \mathtt{7F} \]
The leading \bytes{F0}, the checksum byte itself and the trailing \bytes{F7} take no part
in the sum. Equivalently, the sum of every byte after \bytes{F0} including the checksum is
zero modulo 128.

\begin{caution}{The checksum is not verified for every family}
On reception the checksum is checked only for families \bytes{7E}, \bytes{2B},
\bytes{2C} and \bytes{2D}. For the other accepted families a wrong checksum passes
unremarked.
\end{caution}

\section{Requesting a dump}
A dump can be asked for over MIDI; the front-panel \textsc{send} button is not the only
way to start one. The request names a category by its region byte:

\begin{center}
\bytes{F0 50 2B 04 \textit{nn} 11 \textit{area} 00 00 \textit{x x x} F7}
\end{center}

\begin{center}
\begin{tabular}{ll}
\toprule
\bytes{\textit{area}} & category requested \\
\midrule
\bytes{40} & \textsc{system, part \& midi} \\
\bytes{20} & \textsc{sound} \\
\bytes{50} & \textsc{combination} \\
\bytes{60} & \textsc{sequencer} \\
\bottomrule
\end{tabular}
\end{center}

The three bytes in the length position are ignored; any values are accepted. The
instrument then runs exactly the transfer the corresponding menu row would have run.

There is no request for \textsc{total keyboard}: the four categories must be asked for one
at a time.

\section{Opening a transfer}
A dump begins with a two-step exchange. Each step is attempted up to three times.

\begin{enumerate}
\item The sender transmits \bytes{F0 50 21 04 00 11 F7} and expects
  \bytes{F0 50 22 04 \textit{nn} 11 F7}.
\item The sender transmits \bytes{F0 50 22 04 00 11 F7} and expects
  \bytes{F0 50 23 7E F7}.
\end{enumerate}

If the exchange succeeds, every subsequent data message must be acknowledged. If all three
attempts of the first step fail, \textbf{the dump proceeds anyway, without
acknowledgements}, at a fixed pace of roughly 51\,ms between messages.

\section{Acknowledgement and retry}
With acknowledgements in force, each data message must be answered with
\bytes{F0 50 23 7E F7}. Any other reply causes the same message to be retransmitted, up to
three attempts, after which the transfer is abandoned.

A reply is awaited for about 5.1 seconds; during the opening exchange the wait is about
2.05 seconds.

\section{Continuation}
Data messages after the first open with \bytes{F0 50 7E}, followed by payload,
continuation flag, checksum and \bytes{F7}.

\section{Closing}
At the end of each category the sender transmits \bytes{F0 50 27 7E F7} and waits for an
acknowledgement. At the end of the whole transfer it transmits \bytes{F0 50 28 7E F7}. If
the transfer was abandoned after three failed retries it transmits \bytes{F0 50 29 7E F7}.

\section{Receiving}
When the SX-WSA1R is the receiver it answers each data message with
\bytes{F0 50 23 7E F7} to accept, \bytes{F0 50 24 7E F7} to reject, and
\bytes{F0 50 2A 7E F7} when the incoming data is larger than the destination.

While a bulk-dump session is open the instrument accepts only the \bytes{50} identifier;
universal messages are not honoured until the session ends.

\begin{caution}{The instrument must be on the SYSEX BULK DUMP screen to receive}
A dump cannot be pushed to an idle instrument. Reception happens only while that screen is
displayed; a dump sent at any other time is parsed, discarded, and not reported.

A dump offered on the MIDI~2 input is never received either, at any time. This is what lies
behind the published instruction to use MIDI~1 for bulk transfers.
\end{caution}

Seven of the eight patterns are pinned on all six bytes, address and length together.
Sending a correct address with a shortened length is not a short transfer --- it is an
unrecognised message, refused with \textsc{error 41}. Only \textsc{sequencer} block~3
has its length read as a number, because that block's size varies. An arbitrary
destination cannot be written.

The one length the message itself controls is \textsc{sequencer} part~3, which is capped at
330\,752 bytes; beyond that, and whenever the destination fills mid-message, the
instrument answers \bytes{F0 50 2A 7E F7} and reports \textsc{error 21}.

\begin{caution}{Memory-full is treated as end-of-dump}
When the SX-WSA1R is \emph{sending} and the receiver answers \bytes{F0 50 2A 7E F7} to
report that its memory is full, the instrument treats that exactly as it treats
\bytes{F0 50 28 7E F7}, end of dump. It stops, and it reports success. A sender that needs
to know whether the data arrived cannot rely on the instrument to notice.
\end{caution}
